\chapter{緒言}\label{ch:introduction}
\section{研究背景}
\subsection{沸騰開始を音響から捉える意義}
沸騰を伴う熱流動系では，熱流束の変化とともに伝熱面上の気泡発生状態が変わる．核沸騰開始（onset of nucleate boiling：ONB）付近の状態を捉えることは，沸騰への移行を監視するための基礎的な課題である．本研究の応用背景は，ナトリウム冷却高速炉の炉心局所異常に伴う冷却材沸騰を早期に把握する音響監視技術である．この技術には，沸騰が十分に発達した状態を識別するだけでなく，背景音が存在する条件で開始付近の変化を捉えることが求められる．実際，常陽で行われた音響法による沸騰検出の予備試験では，ポンプ振動やナトリウム流動などによる背景雑音のため，音響信号による気泡発生の明瞭な識別が困難だったことが報告されている\cite{pnc1991}．したがって，雑音下の見逃し・誤報と判断根拠を扱うことは，音響監視技術へ近づくための研究課題となる．

本論文は，この応用に向けた基礎研究として，水を用いたプール沸騰の既存音響データを対象とする．音響から連続的な熱流束を推定することは，熱負荷を定量的に把握する課題であり，ONBを基準とした判定は沸騰への移行を把握する課題である．両者を同じ予測から検討し，推定精度の改善が判定のどの範囲へ及ぶかを明らかにする．水とナトリウムでは物性，伝熱条件，音響伝播および背景音が異なるため，水の試験で得た推定精度を実機の検知性能へ直接読み替えることはできない．水のONB段階の判定結果から，実機での警報時間や事故防止効果を直接示すものではない．

音響計測は，気泡の発生に伴う圧力変動を観測できる一方，計測値には周囲の流動，装置の振動，伝達経路および収録系の影響も含まれる．特定の周波数に強い成分が現れたとしても，それだけで気泡の生成・凝縮や加熱体の共振を一意に同定できるわけではない．手塚の実験研究\cite{tezuka2026}は，伝熱面形状，サブクール度，音響・画像同期計測を比較し，音源の解釈には条件を変えた実験的根拠が必要であることを示している．学習モデルがある帯域を利用することと，その帯域の音源が物理的に確定していることは，区別する必要がある．

また，沸騰が十分に発達した領域では，熱流束の違いに対応した音響変化を学習できても，ONB直後の弱い音響変化を安定して検知できるとは限らない．全域の回帰誤差が小さくても，閾値近傍の過小予測が見逃しを生む可能性がある．したがって，早期検知への適合性を評価するには，熱流束推定の精度とONB付近の判定を，同じ予測値から別々に検討する必要がある．

\subsection{雑音と統合に関する課題}
音響に背景音が重なると，学習時と異なる入力分布が生じる．異なるモデルの予測を統合すれば誤差を抑えられる可能性があるが，モデルが同じ方向に大きく外れる場合には，単に平均するだけで誤差が解消するとは限らない．また，無雑音で良いモデルへ大きい重みを与えることが，雑音下でも有効とは限らない．統合の有効性は，個々のモデルの精度，誤差の方向・大きさ，入力に対する依存，重みの決め方に関係する．

雑音耐性の評価にも二つの問いがある．一つは，雑音を含む条件でどの程度小さい誤差が得られるかである．もう一つは，無雑音からどの程度悪化するかである．無雑音の精度が異なる方式では，雑音下の絶対誤差が小さいことと，悪化量が小さいことは一致しない．本研究では両方を示し，「雑音耐性」という語に異なる意味が混在しないようにする．

\section{音響による熱流束推定・沸騰検知に関する従来研究}
\subsection{音響による定量推定と状態判定}
Dunlapら\cite{dunlap2023}は，ハイドロフォンで取得したプール沸騰音響を用い，特徴抽出と回帰器を組み合わせて熱流束を推定した．出版社掲載の要旨とHighlightsでは，音響区間長と周波数条件が予測に関係することが報告されている．この研究により，音響から熱流束を回帰する枠組み自体には先行研究があることが分かる．ただし，その結果を本研究と同じONB判定，録音分割，付加雑音条件で得た成績として比較することはできない．本稿では数値の直接対決を行わない．

Barathulaら\cite{barathula2023}は，プール沸騰音響から未沸騰，核沸騰，限界熱流束に関係する状態を機械学習で分類している．従って，「音響を使った沸騰検知」や「木モデルの利用」だけを本研究の新規性としては置かない．回帰を介した判定では連続量の誤差と閾値判定の誤りを同じ出力から調べられるが，これだけで直接分類に対する優位が証明されるわけではない．

Uekiら\cite{ueki2025}は，水のサブクール沸騰音響にCNNを適用し，状態分類，Grad-CAMによる利用帯域の可視化，熱流束回帰を報告している．また，植木\cite{ueki2026}は，高速炉での沸騰把握という応用目標と，水を用いる音響検知の基礎検討を位置づけている．これらを踏まえると，回帰，判定および説明図を用いることに加え，モデル統合によって何が改善し，開始段階で何が残るかを明示する必要がある．

Huangら\cite{huang2026}は，サブクール流動沸騰の音響をSTFTへ変換してCNN回帰を行い，白色雑音と運転条件変化の影響を検討している．従って，雑音下の音響回帰という課題自体にも先行研究がある．本研究は，プール沸騰の既存録音に対する固定基準の水流音負荷とモデル統合を扱う．対象，雑音種およびSNRの基準が異なるため，論文間の誤差やSNR値の大小をそのまま性能優位の根拠とはしない．

\subsection{特徴利用の説明と未解決の位置づけ}
回帰の精度が向上した理由を調べるには，モデルがどの入力に依存し，その依存が誤差の変化とどのように対応するかを確認する必要がある．勾配に基づくIntegrated Gradients（IG）\cite{sundararajan2017}は，入力と基準入力の間の出力差を特徴へ配分する方法である．しかし，寄与画像に強い領域が現れることは，その領域を除去したときの性能変化や，音源の物理的同定と同じ証拠ではない．本研究では，共通の入力加工と特徴除去後の学習も用いて，説明画像から言える範囲を補う．

以上を踏まえ，本論文が取り組む課題は，異なる表現・回帰器の統合による全域熱流束推定とONB判定の利得をそれぞれ確かめ，その成立範囲と残る誤りを，誤差構造，入力依存および統合制約から説明することである．従来研究にONB評価が存在しないと包括的に主張するのではなく，本研究のデータ・判定・雑音評価の範囲を具体的に定めて位置づける．
\need{L01}{Dunlapら・Barathulaらは出版社の要旨・公開抜粋，Uekiら2025と植木2026は公開抄録・書誌までの確認である．これらの本文に基づく実験数，分割，録音長およびONB評価の定義を補う．Huangら2026は公開本文の回帰・分割・雑音評価を確認したが，本研究と同一条件の優位性比較ではない．}

\section{本研究の目的}\label{sec:purpose}
本研究は，音響による熱負荷と沸騰への移行の把握に向け，水のプール沸騰の既存音響データを対象に，背景水流音下での熱流束推定誤差と，予測熱流束に基づくONB判定の見逃し・誤報を低減するためのモデル構成および統合条件を明らかにすることを目的とする．熱流束全域の定量推定とONB判定を二つの主要な性能課題とし，改善と悪化の理由を説明することによって，音響監視手法の設計と適用判断に必要な知見を得る．

このため，入力表現・学習方式の異なる回帰器を比較・統合し，同じ評価区間に対する全域回帰誤差と閾値判定を区別して検証する．判定では総見逃し・誤報に加えて，開始段階とその近傍を評価する．さらに，予測誤差の関係，共通入力加工，特徴除去後の学習および統合制約を用いて，どこで利得が生じ，どの誤りが残るかを分析する．性能評価と要因分析は，この目的へ回答するための手段として位置づける．

この目的に対し，次の三つの問いを設定する．
\begin{enumerate}
\item\label{rq:regression} 異なる入力表現・学習方式のモデルの統合は，対象の測定熱流束範囲と付加水流音条件において，単体・基準構成に比べて全域の熱流束推定誤差をどこまで低減するか．
\item\label{rq:onb} 同じ回帰予測の統合は，既知の日別ONB熱流束に基づく総見逃し・誤報をどう変え，その改善は記録上の開始段階と検知到達熱流束まで及ぶか．
\item\label{rq:explanation} 回帰・判定の改善と悪化を，予測誤差の方向・大きさ，入力特徴への依存および統合制約からどこまで説明でき，物理的解釈に何が残るか．
\end{enumerate}

問い\ref{rq:regression}では全域推定の誤差低減を，問い\ref{rq:onb}では総判定と開始段階の利得・限界を独立に回答する．雑音時の絶対誤差と無雑音からの悪化量も区別する．問い\ref{rq:explanation}では，誤差の代数的な説明，モデルの入力依存，物理的な機構を段階に分けて考察する．異なる依存が観測されたことから，雑音下の利得の物理的原因が確定したとはしない．

本論文における「開始に近い検知」は記録上の熱流束段階について検討し，真の気泡発生時刻に対する検知遅延とは区別する．既知録音内の比較という範囲を明記したうえで，全域推定と総判定の成果，開始段階・劣化量・物理的対応に残る限界を結言で回答する．

\begin{table}[htbp]
\centering\small
\caption{研究上の問いと本文での評価・回答の対応．}\label{tab:research_questions}
\begin{tabularx}{\textwidth}{L{13mm}YY}
\toprule
問い & 評価・分析 & 主な回答の配置 \\
\midrule
\ref{rq:regression} & 全域RMSE，雑音時の絶対誤差とclean差，領域別誤差 & 第4章の回帰・雑音評価，第6章 \\
\ref{rq:onb} & 総FN/FP，開始段階，近傍，日別$q_{100}$ & 第4章の判定・近傍評価，第6章 \\
\ref{rq:explanation} & 誤差分解，入力加工，特徴除去，配分・上限診断 & 第5章の要因分析，第6章 \\
\bottomrule
\end{tabularx}
\end{table}

\section{本論文の構成}
第2章では，使用データの出典，熱流束・ONBラベルおよび解析対象を示す．第3章では，音響入力，回帰モデル，学習・統合，雑音の付加，ONB判定および評価指標を定義する．第4章では，統合によって何が改善したかを，全域回帰からONB判定へ段階的に示す．第5章では，第4章の差を領域別誤差，予測の補完および入力依存から考察し，物理的な意味づけの限界を述べる．第6章では研究目的への回答と今後の展望を示す．計算条件，補助結果，未確認事項は付録にまとめる．
